
\section{Názvosloví}

\begin{itemize}
	\item Periodická tabulka prvků -- \textbf{promítnout}, základní orientace
	\item Číselné předpony
	\item Oxidační číslo -- definice, ukázka
	\item Koncovky oxidačních čísel
\end{itemize}

\begin{tabular}{|l|l|l|l|}
	\hline
	Bohrium, Bh & Curium, Cm & Darmstadtium, Ds & Einsteinium, Es \\\hline
	Flerovium, Fl & Hassium, Hs & Kalifornium, Cf & Kopernicium, Cn \\\hline
	Livermorium, Lv & Lutecium, Lu & Meitnerium, Mt & Promethium, Pm \\\hline
	Rhenium, Re & Rhodium, Rh & Roentgenium, Rg & Ruthenium, Ru \\\hline
	Rutherfordium, Rf & Seaborgium, Sg & Tellur, Te & Thallium, Tl \\\hline
	Thulium, Tm & Ytterbium, Yb & Yttrium, Y &  \\\hline
\end{tabular}

\begin{tabular}{|l|l|l|l|l|l|}
	\hline
	$\frac{1}{2}$ & hemi & 5 & penta & 10 & deka \\\hline
	1 & mono & 6 & hexa & 11 & undeka \\\hline
	2 & di & 7 & hepta & 12 & dodeka \\\hline
	3 & tri & 8 & okta & 13 & trideka \\\hline
	4 & tetra & 9 & nona & 14 & tetradeka \\\hline
\end{tabular}

Modrá skalice, zelená skalice, sádra

\begin{tabular}{|c|l|}
	\hline
	\textbf{I} & HCl, HBr, HClO, NaCl, KBr, \ce{K2SO4} \\\hline
	\hline
	\textbf{II} & \ce{H2O}, \ce{MnCl2}, \ce{FeBr2}, \ce{CaSO4} \\\hline
	\hline
	\textbf{III} & \ce{NH3}, \ce{AsH3}, \ce{SbH3}, \ce{HClO2}, \ce{MnF3}, \ce{Fe2O3}, \ce{Al2(SO4)3} \\\hline
	\hline
	\textbf{IV} & \ce{SiH4}, \ce{GeH4}, \ce{H2CO3}, \ce{GeO2}, \ce{GeCl4}, \ce{SnBr4}, \ce{Mn(SO4)2} \\\hline
	\hline
	\textbf{V} & \ce{PH5}, \ce{SbF5}, \ce{HNO3}, \ce{H3PO4}, \ce{N2O5}, \ce{PCl5}, \ce{AsBr5} \\\hline
	\hline
	\textbf{VI} & \ce{SeO3}, \ce{CrO3}, \ce{H2CrO4}, \ce{H6TeO6}, \ce{Na2CrO4}, \ce{K2FeO4}, \ce{WF6} \\\hline
	\hline
	\textbf{VII} & \ce{IF7}, \ce{Mn2O7}, \ce{HClO4}, \ce{H5IO6}, \ce{KMnO4}, \ce{KClO4}, \ce{NaTcO4} \\\hline
	\hline
	\textbf{VIII} & \ce{OsO4}, \ce{H2OsO5} \\\hline
	\hline
	\textbf{0} & \ce{Fe(CO)5}, \ce{Ni(CO)4} \\\hline
\end{tabular}

Určování oxidačního čísla: \ce{O2}, \ce{H2O}, \ce{NH3}, \ce{CaH2}, \ce{K2O}, \ce{K2O2}, \ce{KO2}, \ce{KO3}, \ce{HCl}, \ce{HNO3}, \ce{H2SO4}, \ce{H3PO4}, \ce{H7IO7}

Příklady kationtů: \ce{Na+}, \ce{Al^{3+}}, \ce{Mn^{4+}}, \ce{Hg^{2+}}, \ce{Hg$_2^{2+}$}

Příklady aniontů: \ce{Cl-}, \ce{ClO-}, \ce{ClO$_3^{-}$}

Hydridy, hydroxidy, oxidy, peroxidy, superoxidy, ozonidy, sulfidy, disulfidy, sírany, hydrogensírany, kyseliny fosforu, dusičnany, dusitany

Karbidy: acetylidy (\ce{CaC2}), odvozené od methanu (\ce{WC})

Kyseliny a soli (neutralizace), hydrogensoli

Kyseliny síry: sírová, siřičitá, disírová, kyselina thiosírová, dithionan, tetrathionan

Kyseliny fosforu: metafosforečná -- \ce{(HPO3)_n}, trihydrogenfosforečná, pyrofosforečná (difosforečná), fosforitá, fosforná a jejich soli a organické deriváty -- fosfonáty a fosfináty.

Ionty odvozené od amoniaku: \ce{NH3}, \ce{NH$_4^+$}, \ce{NH$_2^-$}, \ce{NH$^{2-}$}, \ce{N$^{3-}$} a azoimidu, \ce{N$_3^-$}.

Pyridinium, anilinium, acetacidium, methylammonium, fenylammonium.

\subsection{Podvojné soli}

Ve vzorcích podvojných solí se kationty uvádějí v pořadí rostoucích oxidačních čísel, v případě stejného oxidačního čísla v abecedním pořadí. \textit{Víceatomové kationty}, např. \ce{NH$_4^+$} nebo \ce{PH$_4^+$} se uvádějí poslední ve skupině kationtů stejného náboje. V názvu se oddělují pomlčkou a pořadí je dáno pořadím ve vzorci.

\begin{tabular}{|l|l|}
	\hline
	\ce{K2NH4PO4} & fosforečnan didraselno-amonný \\\hline
	\ce{CaMg(SO4)2} & síran vápenato-hořečnatý \\\hline
	\ce{NH4MgPO4} & fosforečnan hořečnato-amonný \\\hline
	\ce{Na(NH4)SO4} & síran sodno-amonný \\\hline
	\ce{KAl(SO4)2.12H2O} & síran draselno-hlinitý \\\hline
	\ce{NaNbO3} & oxid sodno-niobičný (trioxid) \\\hline
\end{tabular}

Anionty se uvádějí v abecedním pořadí značek prvků, příp. centrálních atomů. V případě víceprvkových aniontů se používají číselné předpony bis, tris, …

\begin{tabular}{|l|l|}
	\hline
	\ce{BiCl(O)} & chlorid-oxid bismutitý \\\hline
	\ce{AlO(OH))} & oxid-hydroxid hlinitý \\\hline
	\ce{Ca5F(PO4)3} & fluorid-trisfosforečnan pentavápentý \\\hline
	\ce{Na6ClF(SO4)2} & chlorid-fluorid-bis(síran) hexasodný \\\hline
	\ce{Na4ClF(S2O7)} & chlorid-fluorid-disíran tetrasodný \\\hline
	\ce{Ca3(CO3)2F2} & bis(uhličitan)-difluorid trivápenatý \\\hline
\end{tabular}

\subsection{Iso- a heteropolyionty}

Síran, disíran, disiřičitan!, fosforečnan, difosforečnan (\textit{odvození pomocí dehydratace}), \textit{katena}-trifosforečnan, \textit{cyklo}-trifosforečnan, chroman, dichroman.

\subsection{Funkční deriváty kyseliny}

Vznik, halogenidy, amidy a estery

Kationtové skupiny:

\begin{tabular}{|l|l|l|l|}
	\hline
	OH & hydroxyl & SeO & seleninyl \\\hline
	CO & karbonyl & \ce{SeO2} & selenonyl \\\hline
	NO & nitrosyl & \ce{CrO2} & chromyl \\\hline
	\ce{NO2} & nitryl & \ce{UO2} & uranyl \\\hline
	PO & fosforyl & ClO & chlorosyl \\\hline
	VO & vanadyl & \ce{ClO2} & chloryl \\\hline
	SO & thionyl & \ce{ClO3} & perchloryl \\\hline
	\ce{SO2} & sulfuryl & \ce{S2O5} & disulfuryl \\\hline
\end{tabular}

Thiokyseliny, dithionan, trithionan, tetrathionan.

\ce{SO2(NH2)2}, \ce{SO(NH2)}, \ce{CO(NH2)2} (močovina), \ce{OP(OEt)3}, ethylester kyseliny octové

\subsection{Názvoslovná rozcvička}

\begin{tabular}{|l||l|}
	\hline
	\ce{Al2S3} & disulfid sodný \\\hline
	\ce{Mo(SO4)2} & kyanid draselný \\\hline
	\ce{NaNO2} & thiokyanatan amonný \\\hline
	\ce{NaI3} & thiosíran draselný \\\hline
	\ce{CaO2} & oxid selenový \\\hline
	\ce{NaO2} & oxid osmičelý \\\hline
	\ce{NaClO} & dodekahydrát síranu draselno-hlinitého \\\hline
	\ce{RbClO2} & bromitan vápenatý \\\hline
	\ce{KBrO3} & trimethylester kyseliny borité \\\hline
	\ce{CsClO4} & oxid křemičitý \\\hline
	\ce{K2Cr2O7} & křemičitan sodný \\\hline
	\ce{CrPO4} & arseničnan hořečnatý \\\hline
	\ce{K5P3O10} & fluorid jodistý \\\hline
	\ce{K3P3O9} & tetrathionan draselný \\\hline
	\ce{Na2H3P3O10} & disíran barnatý \\\hline
	\ce{CuSO4.5H2O} & síran tetrafenylfosfonia \\\hline
	\ce{(NEt4)2SO4} & alan \\\hline
	\ce{KNH2} & wolframan železnatý \\\hline
	\ce{H5IO6} & peroxid rubidný \\\hline
	\ce{KSO3F} & kyanatan sodný \\\hline
	\ce{COCl2} & hydrogenuhličitan vápenatý \\\hline
	\ce{KOCN} & trichlorid vanadylu \\\hline
	\ce{Na2S2O3} & chlorid měďný \\\hline
	\ce{Hg2Cl2} & hydrogensulfid hlinitý \\\hline
\end{tabular}

\clearpage